\cleardoublepage
\phantomsection
\addcontentsline{toc}{chapter}{Avant-propos}
\chapter*{Avant-propos}



Pour faciliter l'intégration des étudiants dans le milieu professionnel, l'École Nationale de Commerce et de Gestion d'Agadir (ENCGA) encourage ses étudiants à effectuer un stage en entreprise. Cette expérience leur permet de confronter la théorie acquise en formation à la réalité du terrain.

C'est dans ce cadre que s'inscrit mon stage de perfectionnement au sein du Centre d'Affaires de Bank of Africa, réalisé du 01/05/2026 au 30/06/2026.

Le choix de la thématique « Analyse Comparative Entre Les Modèles De Rating De Risque De Crédit » résulte d'une décision réfléchie et planifiée. La gestion des risques constituant une composante majeure de la finance quantitative, ce sujet est en parfaite cohérence avec mon projet professionnel, qui est d'évoluer dans le domaine de l'ingénierie et de la finance quantitative. Il correspond pleinement à mon intérêt pour l'application des méthodes quantitatives aux problématiques financières.

À travers cette mission, j'ai cherché à mobiliser mes connaissances, à observer de manière active, à analyser les problématiques rencontrées et à proposer des pistes d'amélioration. Ce rapport reflète, je l'espère, non seulement la richesse de cette expérience professionnelle, mais aussi l'implication personnelle que j'y ai consacrée.